
\documentclass[../../main.tex]{subfiles}
\begin{document}

% 年度の大きな表示
\begin{center}
  {\fontsize{48pt}{58pt}\selectfont \textbf{1996}\normalsize\textbf{年度}}
\end{center}

\vspace{10mm}

% 本文


\vspace{15mm}

% 出題テーマと難易度の表
\noindent\textbf{出題テーマと難易度}

\vspace{5mm}

\begin{center}
  \shadowbox{%
    \begin{tabular}{|c|p{8cm}|c|}
      \hline
      \textbf{問題番号} & \textbf{テーマ} & \textbf{難易度} \\
      \hline
      第1問 & 整数問題（連立方程式の正整数解） & \\
      \hline
      第2問 & 行列と一次変換（楕円上の像） & \\
      \hline
      第3問 & 高次関数の極値と面積条件 & \\
      \hline
      第4問 & 微分方程式的関係と接線・法線の性質 & \\
      \hline
    \end{tabular}%
  }
\end{center}

\vspace{15mm}

% 解答
\noindent\textbf{解答}

\vspace{5mm}

\begin{center}
  \shadowbox{%
    \renewcommand{\arraystretch}{1.6}
    \begin{tabular}{|c|c|p{8cm}|}
      \hline
      \textbf{問題番号} & \textbf{小問} & \textbf{答え} \\
      \hline
      \multirow{3}{*}{第1問} & (1) & $(x_1,x_2)=(2,2)$（$n=2$），$(x_1,x_2,x_3)$は$(1,2,3)$の並べ替え6通り（$n=3$） \\
      \cline{2-3}
                             & (2) & $n=2$ \\
      \cline{2-3}
                             & (3) & 証明問題 \\
      \hline
      第2問 & - & （要確認） \\
      \hline
      \multirow{2}{*}{第3問} & (1) & $f(x)=10x^7\left(x-\dfrac{6}{5}\right)\left(x-\dfrac{3}{2}\right)$，$f(x)=-98x^7\left(x-\dfrac{6}{7}\right)\left(x-\dfrac{15}{14}\right)$ \\
      \cline{2-3}
                             & (2) & $f(x)=-\dfrac{7}{2}x^7\left(x^2-\dfrac{9}{7}\right)$ \\
      \hline
      \multirow{2}{*}{第4問} & (1) & 証明問題 \\
      \cline{2-3}
                             & (2) & $\min F(t)=\dfrac{16}{9}\sqrt{3}$（$1<a\le\dfrac{4}{3}$），$\dfrac{a^2}{\sqrt{a-1}}$（$a\ge\dfrac{4}{3}$） \\
      \hline
    \end{tabular}%
    \renewcommand{\arraystretch}{1.0}
  }
\end{center}

\vspace{15mm}

% 解答時間
\noindent\textbf{試験形態}

\vspace{5mm}

\begin{center}
  \shadowbox{%
    \begin{tabular}{p{8cm}}
      （要確認）
    \end{tabular}%
  }
\end{center}

\end{document}
